
\documentclass[addpoints]{exam}
\usepackage{amsmath}
\usepackage{amssymb}
\usepackage{color}

% \printanswers

\newcommand{\event}{Geometric and Cartesian Vectors}
\newcommand{\tournament}{}
\def\type{0}

\setlength\fillinlinelength{\textwidth}
\CorrectChoiceEmphasis{\color{red}\bfseries}
\SolutionEmphasis{\color{red}\bfseries}

\setlength\answerclearance{2pt}
\pagestyle{headandfoot}
\runningheadrule
\runningheader{\event}{\tournament}{\ifnum\type=1 Class Set Test \else Full Name:\hspace{1cm} \fi}
\runningfootrule
\runningfooter{}{Page \thepage\ of \numpages}{}

\begin{document}
\begin{center}
    {\huge Grade 12 Calculus \& Vectors \\ \event
    \ifprintanswers \space Key
    \else \space \ifnum\type<2 Test \fi \ifnum\type=1 \textbf{(CLASS SET)} \fi
          \ifnum\type=2 Answer Sheet \fi
    \fi \\ \vspace{3mm}\tournament\vspace{1mm}}
\end{center}

\vspace{.25\textheight}

\begin{center}
    First Name: \ifprintanswers
    \underline{\hspace{3.5cm}}\textbf{KEY}\underline{\hspace{5.5cm}}\\\vspace{5mm}
    \else \underline{\hspace{10cm}}\\ \vspace{5mm} \fi
    Last Name: \underline{\hspace{10cm}}\\ \vspace{5mm}
\end{center}

\noindent\textbf{\underline{Directions:}}
\begin{itemize}
    \ifnum\type=1 \item \textbf{This test is a class set.} Do not write on this paper. \fi
    \item Show all work. Express exact answers (leave surds unsimplified unless asked).
    \item The test is designed to be completed in 75 minutes.
\end{itemize}
\begin{center}
\textit{For grading use only}\\
\multirowgradetable{1}[pages]
\end{center}

\newpage
\begin{questions}

\section*{Part A --- Multiple Choice (15 marks)}
\vspace{5mm}

\question[1] A vector is completely described by its:
\begin{choices}
    \choice Magnitude only
    \choice Direction only
    \correctchoice Both magnitude and direction
    \choice Position in the plane
\end{choices}

\question[1] If $\vec{a} = (3,\, 4)$, then $|\vec{a}|$ equals:
\begin{choices}
    \choice 3.5
    \correctchoice 5
    \choice 7
    \choice 25
\end{choices}

\question[1] If $\vec{a} \perp \vec{b}$, then $\vec{a} \cdot \vec{b}$ equals:
\begin{choices}
    \choice 1
    \choice $|\vec{a}||\vec{b}|$
    \correctchoice 0
    \choice $-1$
\end{choices}

\question[1] For $\vec{a} = (1,\, 2,\, -1)$ and $\vec{b} = (3,\, -1,\, 2)$,
$\vec{a} \cdot \vec{b}$ equals:
\begin{choices}
    \correctchoice $-1$
    \choice $1$
    \choice $3$
    \choice $-3$
\end{choices}

\question[1] The cross product $\vec{a} \times \vec{b}$ is:
\begin{choices}
    \choice A scalar
    \choice Always zero
    \correctchoice A vector perpendicular to both $\vec{a}$ and $\vec{b}$
    \choice Parallel to $\vec{a}$
\end{choices}

\question[1] $|\vec{a} \times \vec{b}|$ equals the area of:
\begin{choices}
    \choice The triangle with sides $\vec{a}$ and $\vec{b}$
    \correctchoice The parallelogram with sides $\vec{a}$ and $\vec{b}$
    \choice A circle of radius $|\vec{a}|$
    \choice A trapezoid
\end{choices}

\question[1] The unit vector in the direction of $\vec{a} = (0,\, 3,\, 4)$ is:
\begin{choices}
    \choice $(0,\, 1/3,\, 1/4)$
    \correctchoice $(0,\, 3/5,\, 4/5)$
    \choice $(0,\, 9,\, 16)$
    \choice $(0,\, 3,\, 4)/25$
\end{choices}

\question[1] The angle between $\vec{a} = (1,\, 0)$ and $\vec{b} = (1,\, 1)$ is:
\begin{choices}
    \choice $30^\circ$
    \correctchoice $45^\circ$
    \choice $60^\circ$
    \choice $90^\circ$
\end{choices}

\question[1] The scalar projection of $\vec{a}$ onto $\vec{b}$ is:
\begin{choices}
    \choice $\vec{a} \cdot \vec{b}$
    \choice $\dfrac{\vec{a} \cdot \vec{b}}{|\vec{b}|^2}$
    \correctchoice $\dfrac{\vec{a} \cdot \vec{b}}{|\vec{b}|}$
    \choice $|\vec{a}|\sin\theta$
\end{choices}

\question[1] $\vec{i} \times \vec{j}$ equals:
\begin{choices}
    \choice $\vec{i}$
    \choice $\vec{j}$
    \correctchoice $\vec{k}$
    \choice $-\vec{k}$
\end{choices}

\question[1] The work done by force $\vec{F} = (3,\, 4)$ over displacement $\vec{d} = (2,\, 1)$
is:
\begin{choices}
    \choice 5
    \correctchoice 10
    \choice 14
    \choice 20
\end{choices}

\question[1] For $\vec{a} = (2,\, -1,\, 3)$ and $\vec{b} = (k,\, 2,\, -1)$, if
$\vec{a} \perp \vec{b}$, then $k = $
\begin{choices}
    \choice $1$
    \correctchoice $5/2$
    \choice $3$
    \choice $-5/2$
\end{choices}

\question[1] Vectors $\vec{a} = (1,\, 2)$ and $\vec{b} = (2,\, 4)$ are:
\begin{choices}
    \choice Perpendicular
    \choice Equal
    \correctchoice Parallel (collinear)
    \choice Skew
\end{choices}

\question[1] $|\vec{a} \times \vec{b}|^2 + (\vec{a} \cdot \vec{b})^2 = $
\begin{choices}
    \choice $0$
    \choice $|\vec{a}|^2 + |\vec{b}|^2$
    \correctchoice $|\vec{a}|^2|\vec{b}|^2$
    \choice $|\vec{a}||\vec{b}|$
\end{choices}

\question[1] If $\vec{a} = (1,\, 2,\, 3)$, $\vec{b} = (4,\, 5,\, 6)$, then
$2\vec{a} - \vec{b}$ equals:
\begin{choices}
    \correctchoice $(-2,\, -1,\, 0)$
    \choice $(2,\, 1,\, 0)$
    \choice $(6,\, 9,\, 12)$
    \choice $(-2,\, 1,\, 0)$
\end{choices}


\section*{Part B --- Short Answer (33 marks)}

\question Let $\vec{a} = (2,\,-1,\,3)$, $\vec{b} = (-1,\,4,\,2)$, $\vec{c} = (0,\,3,\,-2)$.
\begin{parts}
    \part[2] Find $2\vec{a} - \vec{b}$.
    \part[2] Find $|\vec{a} + \vec{c}|$.
    \part[2] Find the unit vector in the direction of $\vec{b}$.
    \part[2] Find the value of $t$ for which $\vec{a} + t\vec{b}$ is perpendicular to $\vec{c}$.
\end{parts}
\begin{solution}[9cm]
\textbf{(a)} $2\vec{a} - \vec{b} = (4,-2,6) - (-1,4,2) = \mathbf{(5,\,-6,\,4)}$

\textbf{(b)} $\vec{a} + \vec{c} = (2,2,1)$.\quad $|(2,2,1)| = \sqrt{4+4+1} = \mathbf{3}$

\textbf{(c)} $|\vec{b}| = \sqrt{1+16+4} = \sqrt{21}$.\quad
Unit vector $= \dfrac{1}{\sqrt{21}}(-1,4,2)$

\textbf{(d)} $(\vec{a}+t\vec{b})\cdot\vec{c} = 0$:
$(2-t,\,-1+4t,\,3+2t)\cdot(0,3,-2) = 0$
\[
3(-1+4t) + (-2)(3+2t) = -3+12t-6-4t = -9+8t = 0 \implies \boxed{t = \frac{9}{8}}
\]
\end{solution}

\question Let $\vec{u} = (1,\,3,\,-2)$ and $\vec{v} = (2,\,-1,\,4)$.
\begin{parts}
    \part[2] Calculate $\vec{u} \cdot \vec{v}$.
    \part[2] Find the angle between $\vec{u}$ and $\vec{v}$, to the nearest degree.
    \part[2] Find the scalar projection of $\vec{u}$ onto $\vec{v}$.
    \part[2] Find the vector projection of $\vec{u}$ onto $\vec{v}$.
\end{parts}
\begin{solution}[9cm]
\textbf{(a)} $\vec{u}\cdot\vec{v} = 2 - 3 - 8 = \mathbf{-9}$

\textbf{(b)} $|\vec{u}| = \sqrt{1+9+4} = \sqrt{14}$;\quad $|\vec{v}| = \sqrt{4+1+16} = \sqrt{21}$
\[
\cos\theta = \frac{-9}{\sqrt{14}\cdot\sqrt{21}} = \frac{-9}{\sqrt{294}} = \frac{-9}{7\sqrt{6}}
\approx -0.5249 \implies \theta \approx \mathbf{122^\circ}
\]

\textbf{(c)} Scalar projection $= \dfrac{\vec{u}\cdot\vec{v}}{|\vec{v}|} = \dfrac{-9}{\sqrt{21}}
= \dfrac{-9\sqrt{21}}{21} = \mathbf{-\dfrac{3\sqrt{21}}{7}}$

\textbf{(d)} Vector projection
$= \dfrac{\vec{u}\cdot\vec{v}}{|\vec{v}|^2}\vec{v} = \dfrac{-9}{21}(2,-1,4)
= -\dfrac{3}{7}(2,-1,4) = \mathbf{\left(-\dfrac{6}{7},\,\dfrac{3}{7},\,-\dfrac{12}{7}\right)}$
\end{solution}

\question Let $\vec{p} = (1,\,2,\,3)$ and $\vec{q} = (2,\,-1,\,1)$.
\begin{parts}
    \part[3] Calculate $\vec{p} \times \vec{q}$.
    \part[2] Verify that $\vec{p} \times \vec{q}$ is perpendicular to both $\vec{p}$ and $\vec{q}$.
    \part[2] Find the area of the parallelogram with adjacent sides $\vec{p}$ and $\vec{q}$.
    \part[1] Find the area of the triangle with these two sides.
\end{parts}
\begin{solution}[9cm]
\textbf{(a)}
\[
\vec{p} \times \vec{q} = \begin{vmatrix}
\vec{i} & \vec{j} & \vec{k} \\
1 & 2 & 3 \\
2 & -1 & 1
\end{vmatrix}
= \vec{i}(2\cdot1 - 3\cdot(-1)) - \vec{j}(1\cdot1 - 3\cdot2) + \vec{k}(1\cdot(-1) - 2\cdot2)
= \mathbf{(5,\,5,\,-5)}
\]

\textbf{(b)}
$(5,5,-5)\cdot(1,2,3) = 5+10-15 = 0$\checkmark;\quad
$(5,5,-5)\cdot(2,-1,1) = 10-5-5 = 0$\checkmark

\textbf{(c)} $|\vec{p}\times\vec{q}| = \sqrt{25+25+25} = \sqrt{75} = \mathbf{5\sqrt{3}}$ units$^2$

\textbf{(d)} Area of triangle $= \dfrac{1}{2}|\vec{p}\times\vec{q}| = \mathbf{\dfrac{5\sqrt{3}}{2}}$ units$^2$
\end{solution}

\question A sled is pulled with a force of 80 N at an angle of $30^\circ$ above the horizontal.
The sled moves 15 m along a horizontal surface.
\begin{parts}
    \part[2] Calculate the work done by the applied force.
    \part[2] If the angle is increased to $45^\circ$ (force unchanged), how does the work
    change? Calculate the new work.
    \part[2] At what angle $\theta$ is the work maximized, and what is that maximum work?
    \part[2] A second force $\vec{F}_2 = (20,\,10,\,0)$ N acts simultaneously. Find the total
    work done by both forces over displacement $\vec{d} = (15,\,0,\,0)$ m. (Use
    $\vec{F}_1 = (80\cos30^\circ,\,80\sin30^\circ,\,0)$.)
\end{parts}
\begin{solution}[9cm]
\textbf{(a)} $W = |\vec{F}||\vec{d}|\cos\theta = 80 \times 15 \times \cos30^\circ
= 1200 \times \dfrac{\sqrt{3}}{2} = \mathbf{600\sqrt{3} \approx 1039 \text{ J}}$

\textbf{(b)} $W = 80 \times 15 \times \cos45^\circ = 1200 \times \dfrac{\sqrt{2}}{2}
= \mathbf{600\sqrt{2} \approx 848 \text{ J}}$
The work \textbf{decreases} as the angle increases.

\textbf{(c)} $W = 80 \times 15 \times \cos\theta$ is maximized when $\cos\theta = 1$,
i.e.\ $\theta = 0^\circ$ (horizontal pull).
$W_{\max} = \mathbf{1200 \text{ J}}$.

\textbf{(d)} $\vec{F}_1 = (40\sqrt{3},\,40,\,0)$.
Total force: $\vec{F}_{total} = (40\sqrt{3}+20,\,50,\,0)$.
\[
W_{total} = \vec{F}_{total}\cdot\vec{d} = (40\sqrt{3}+20)(15) + (50)(0) + 0
= 600\sqrt{3} + 300 \approx \mathbf{1339 \text{ J}}
\]
\end{solution}

\end{questions}
\end{document}
